\section*{Sets and Indices}

\begin{itemize}
    \item $R=\{\mathrm{FT},\mathrm{PT}\}$: clerk types, where $\mathrm{FT}$ denotes full-time and $\mathrm{PT}$ denotes part-time.
\end{itemize}

\section*{Parameters}

\begin{itemize}
    \item $H_r$: regular monthly working hours per clerk of type $r$, from \texttt{Monthly\_Working\_Hours}.
    \item $s_r$: sales volume in pairs per hour for clerk type $r$, from \texttt{Sales\_Volume\_Pairs\_Per\_Hour}.
    \item $N_r$: number of clerks of type $r$. In the implemented code, these default to
    \[
    N_{\mathrm{FT}}=5,\qquad N_{\mathrm{PT}}=4,
    \]
    unless a clerk-count column is detected in the CSV.
    \item $G$ (code: \texttt{sales\_goal}): monthly sales goal.
    \item $\rho$ (code: \texttt{pt\_ft\_ratio}): minimum ratio of part-time overtime to full-time overtime.
    \item $S^{\textnormal{reg}}$: total sales achievable under full regular employment,
    \[
    S^{\textnormal{reg}}=\sum_{r\in R} N_r H_r s_r.
    \]
    \item $D$: sales shortfall to be covered by overtime,
    \[
    D=G-S^{\textnormal{reg}}.
    \]
\end{itemize}

Using the provided data,
\[
H_{\mathrm{FT}}=160,\quad s_{\mathrm{FT}}=5,\quad
H_{\mathrm{PT}}=80,\quad s_{\mathrm{PT}}=2,\quad
G=5500,\quad \rho=0.5.
\]

\section*{Decision Variables}

\begin{itemize}
    \item $y_{\mathrm{FT}}$ (code: \texttt{y\_ft}): overtime hours assigned in total to full-time clerks.
    \item $y_{\mathrm{PT}}$ (code: \texttt{y\_pt}): overtime hours assigned in total to part-time clerks.
\end{itemize}

Both variables are continuous and nonnegative.

\section*{Objective}

Minimize total overtime hours:
\[
\min \; y_{\mathrm{FT}} + y_{\mathrm{PT}}.
\]

\section*{Constraints}

\begin{align}
s_{\mathrm{FT}} y_{\mathrm{FT}} + s_{\mathrm{PT}} y_{\mathrm{PT}} &\ge D
&&\text{(overtime sales must cover the sales shortfall)},\\
y_{\mathrm{PT}} &\ge \rho\, y_{\mathrm{FT}}
&&\text{(part-time overtime coverage requirement)},\\
y_{\mathrm{FT}},\, y_{\mathrm{PT}} &\ge 0.
\end{align}

Equivalently, after substituting $D=G-\sum_{r\in R}N_rH_rs_r$,
\[
s_{\mathrm{FT}} y_{\mathrm{FT}} + s_{\mathrm{PT}} y_{\mathrm{PT}}
\;\ge\;
G-\sum_{r\in R}N_rH_rs_r.
\]

With the default clerk counts used by the code,
\[
S^{\textnormal{reg}}=5\cdot160\cdot5 + 4\cdot80\cdot2 = 4640,
\]
so the implemented sales-shortfall constraint becomes
\[
5y_{\mathrm{FT}} + 2y_{\mathrm{PT}} \ge 5500-4640=860.
\]

\section*{Notes on Implementation}

\begin{itemize}
    \item The code solves a linear program with \texttt{pulp} using \texttt{GUROBI\_CMD}.
    \item Although the business description mentions profit, wages, overtime pay, full employment, and minimizing overtime, the implemented model uses only the sales goal and the overtime-ratio requirement in the optimization. Profit and wage parameters are read from the CSV but do not appear in the mathematical program.
    \item Full employment is not a decision in the code; it is imposed implicitly by computing regular sales from all available clerks and treating only overtime as variable.
    \item The number of clerks is hard-coded as 5 full-time and 4 part-time unless a count-related column is found in the clerk data table.
    \item No upper bounds on overtime are imposed in the implementation.
\end{itemize}
